% !PRINTED 417
\weberpage{417}
%<<CONT
a formula which can easily be generalised by complete induction as follows.

If one denotes by $\mu_1$ all the numbers that arise from $n$ when one divides
$n$ by an even number of distinct prime divisors of $n$ ($n$ itself included),
and by $\mu_2$ the numbers that arise from $n$ when one divides $n$ by an odd
number of such prime divisors, then

\weq{10}{X_n = \dfrac{\prod\limits_{\mu_1}(x^{\mu_1} - 1)}{\prod\limits_{\mu_2}(x^{\mu_2} - 1)}.}

The proof results from (9) for $npp'$, if one assumes that the correctness has
already been proved for $n$.

If $r$ denotes each of the primitive $n$th roots of unity, then for \gesperrt{odd} $n$
the $-r$ are the primitive $2n$th roots of unity (\S.\,132, I, II).
Accordingly, for odd $n$,

\weq{11}{X_{2n}(x) = X_n(-x).}

If $n$ is a power of $2$, then by (8):

\weq{12}{X_n = \dfrac{x^n - 1}{x^{\frac{n}{2}} - 1} = x^{\frac{n}{2}} + 1.}

If we set $x = 1$ in formula (8), then $X_n$ takes the value $p$ if $n$ is a
power of $p$. On the other hand, formula (9) yields, when $n > 1$, the value
$1$ for $X_{npp'}$ (for $n = 1$ the numerator and denominator on the
right-hand side would be $= 0$).

We thus obtain the theorem:

\gesperrt{V. The function $X_n$ takes for $x = 1$ the value $p$ if $n$ is a
power of the prime $p$, and the value $1$ if $n$ has more than one prime as
divisor.}

\begin{center}
\S.\,134.

\gesperrt{Irreducibility.}
\end{center}

The functions $X_n$, which are defined in the preceding section, have the
property of \gesperrt{irreducibility}, which is very important for their
deeper study and all applications, i.e.

\gesperrt{It is not possible to decompose a function $X_n$ of $x$ into factors of lower
degree which themselves}
%CONT>>
